% Interpretation of completed results; no new measurements or labels.
\section{Discussion}\label{sec:discussion}

\paragraph{Morphology (RQ1).}
Corrected lemma expansion increases BM25 nDCG@10 by 0.031 under the primary labels, with a paired 95\% interval from 0.013 to 0.051. This supports a gain for this query set under these assessments. The weaker evidence under Gemma~4 12B limits a judge-independent conclusion. The earlier decrease before the analyzer fix also illustrates an implementation dependency: expansion with an incorrectly guessed lemma can change retrieval in the opposite direction. The result concerns the tested analyzer rule and corpus, rather than morphology as a universally beneficial operation.

\paragraph{Fusion and technical queries (RQ2--RQ3).}
Adaptive fusion exceeds BM25 on the primary endpoint but has no detected advantage over EG2 or RRF. E5 has the highest mean nDCG@10, including in the control with the same text candidates. The absence of a Holm-significant E5--adaptive contrast does not establish equal effectiveness. The phrase-reranking ablation is numerically identical to the full system on the frozen runs; the identifier and short-fragment ablations likewise provide no evidence of an improvement on this sample. A rule can address a plausible failure case without that case occurring often enough here to improve the aggregate metric. A general recommendation about the fusion rule would therefore require independently written queries and broader target coverage.

\paragraph{Changing target coverage.}
The separate extension produces a different ordering: EG2 and adaptive fusion have higher point estimates than E5 under its initial Gemini labels. Query authorship, target formats and the assessment variant also change between cohorts. The extension therefore cannot identify which factor produces that difference, and its means should not be treated as a before/after improvement over the original benchmark. Retaining the 50-record analysis beside the deterministic 49-wording sensitivity preserves the original data while making the duplicated wording visible.

\paragraph{Multimodal matching (RQ4).}
Caption--image and transcript--recording tasks test cross-modal matching in two languages using source-dataset pairings. Image recall at one is 0.857 for Ukrainian captions and 0.889 for Russian captions among 1\,000 images. These outcomes are distinct from the model-judged text-retrieval results. For speech, using the same transcript as query and indexed text is a ceiling condition; even cross-lingual transcripts are parallel source sentences. Neither task establishes effectiveness for a person's paraphrase, a long recording or arbitrary photographs. The tasks support the implemented modality connection within their stated matching conditions.

\paragraph{Quantization and judgement coverage (RQ5).}
The ordinary Q4 comparison has a negative point difference, while its condensed-list comparison does not show the same loss; 9.9\% of Q4 top-10 results are unjudged. The frozen pool therefore cannot identify an intrinsic Q4 relevance loss from the ordinary score. Conversely, the small f16/Q8 differences and intervals spanning zero do not establish equivalence. Backend and precision change together, and the 21 PDF-page rows also change input modality. Quality, file size and encoding cost should be read jointly with these coverage and implementation conditions, rather than as an isolated causal effect of quantization.

\paragraph{The cost of the implemented search path.}
At two million synthetic vectors, the actual archive-ranking path has median time about 2\,442~ms, while the optimized NumPy top-10 reference takes about 107~ms. The paths perform different work: the archive method gathers candidates, fully sorts scores and constructs 400 tuples. This gap identifies implementation work that a matrix-product-only description would miss. HNSW with search depth 256 takes about 1.28~ms and has strict-ID recall@10 of 0.9825 on this synthetic collection. These sequential desktop measurements are not an end-to-end or paired speedup estimate, and the recall need not transfer to a real archive of that size. They motivate approximate indexing and candidate-selection work as future engineering directions.

\paragraph{Assessment and practical interpretation.}
Judge agreement, target checks and deliberate no-answer queries diagnose assessment behaviour without establishing human relevance validity. The Gemini review retains dependence on the original assessors and visible grades, so label sensitivity accompanies the point estimates. An independent human study and queries written by people seeking their own documents remain future validation; they are not reported as completed here. The private-archive checks show that the system can satisfy the author's operational expectations in that deployment, while their development reuse prevents a user-effectiveness claim. Taken together, the evidence supports an inspectable account of quality and cost in one constructed benchmark and one deployment, with the strongest comparisons conditional on their labels, populations and measurement boundaries.
